\documentclass[a4paper]{article}

\usepackage[utf8]{inputenc}
\usepackage[italian]{babel}
\usepackage{graphicx}
\usepackage[a-1a]{pdfx}

\usepackage[a4paper, margin=3cm]{geometry}

\newcommand{\skipline}{\vspace{0.2in}}

\begin{document}

\begin{center}
    \includegraphics[width=0.9\textwidth]{../tesi/Logo.jpg}
\end{center}
\vspace{3mm}
\begin{center}
    {\huge{\bf Studio sull'incidentalità}}\\
    \vspace{2mm}
    {\huge{\bf stradale tramite dataset aperti}}\\
    \vspace{3mm}
    {\bf Gabriele Padovani - Matr. 909165}\\
    \vspace{3mm}
    {\bf Termine: 5 febbraio 2021}\\
\end{center}

\section{Introduzione}

Nel documento di tesi si sono esplorate quali tipi di analisi è possibile realizzare, 
avendo a disposizione una sufficiente quantità di dati liberi. 
Gli argomenti di questo lavoro si concentrano, in particolare, sull'ambito dell'incidentalità 
stradale, dando enfasi alla posizione in cui questi ultimi avvengono, 
sia dal punto di vista delle coordinate geografiche, sia da quello delle strade 
con maggiore numero di sinistri. 

Si è, inoltre, tentato di ricavare il legame tra la quantità di incidenti e svariati altri 
fattori, tra cui, la presenza di passeggeri, l'età del conducente, 
la pavimentazione della direttrice, o la tipologia di incrocio.

\section{Origine dei dati}

Le informazioni utilizzate, generalmente raggruppate in dataset, sono 
state reperite nei principali siti di open data, come quello del comune di 
Milano\footnote{\url{https://dati.comune.milano.it/}} per i dati sugli accessi in area C, 
o quello del ministero dei trasporti\footnote{\url{https://www.mit.gov.it/}}, 
nel caso delle patenti.

D'altra parte, i file principali, contenenti i dati riguardanti gli incidenti, sono 
stati reperiti sul sito 
dell'Istituto nazionale di statistica\footnote{\url{https://www.istat.it/it/archivio/87539}}, 
e su quello dell'Automobile club 
Italia\footnote{\url{http://www.aci.it/fileadmin/documenti/ACI/Trasparenza/Open_Data/Catalogo_localizzazione_in_formato_Open_2019.pdf}}. 

I dataset contenenti le coordinate degli incidenti avvenuti a Milano, 
provengono invece dal giornale online 
TheSubmarine\footnote{\url{https://thesubmarine.it/2018/06/20/mappa-incidenti-stradali-milano}}, 
che ha ottenuto, da parte di Istat, il rilascio di una frazione di questi dati, 
normalmente oscurati per rispettare la privacy degli individui coinvolti nei sinistri. 

Di fondamentale importanza è stato anche il portale di dati geolocalizzati del comune di 
Milano\footnote{\url{https://geoportale.comune.milano.it/}}, 
tramite cui si sono ottenuti i percorsi dei trasporti pubblici, e le 
informazioni sul manto stradale.

Infine, per poter rappresentare i risultati ricavati, in molti casi si è fatto uso di 
alcune mappe, riguardanti le regioni e le province italiane, trovate su un \textit{repository} 
Github\footnote{\url{https://github.com/openpolis/geojson-italy}}.

\section{Analisi}

Nonostante le analisi degli eventi che dispongono di localizzazione 
si siano incentrate soprattutto sulla relazione tra incidentalità e pavimentazione 
della strada, è anche stata posta particolare enfasi sia sul rapporto tra numero di 
incidenti e la presenza di linee di trasporti pubblici, sia, 
allo stesso modo, sulla differenza di sinistri nelle vicinanze di autovelox. 

Questi calcoli prendono in considerazione solamente la città di Milano, 
in quanto le informazioni rilasciate da Istat sono limitate al capoluogo lombardo.

\skipline
Nelle analisi realizzate sui dataset provenienti dal sito Istat, 
ci si è soffermati soprattutto sull'evoluzione dell'incidentalità nei diversi orari 
della giornata, e allo stesso modo al passare dei mesi. 

Si è inoltre controllata l'esistenza di tipologie di strade e incroci che favorissero 
un numero maggiore di sinistri e, allo stesso modo, se alcuni di questi presentassero 
un valore anomalo di pedoni coinvolti. 

Altre analisi riguardanti queste informazioni, si sono concentrate su differenze nella 
guida di uomini e donne, e nelle differenze del numero di sinistri al cambiare della tratta stradale. 

Infine, si è tentato di individuare l'esistenza di fattori che 
potessero distrarre il conducente, in particolare nella forma di telefoni cellulari 
o nella presenza di altri passeggeri a bordo. 

\skipline
I dataset provenienti dal sito Aci, d'altra parte, sono stati oggetto di varie analisi per 
quanto riguarda il luogo dell'incidente, inteso come il nome della strada o la provincia 
in cui questo è avvenuto. 

Più nello specifico, si sono create varie rappresentazioni, incentrate sugli incidenti 
per le diverse regioni italiane, per le province di Lombardia e Lazio, e per le 
principali autostrade, sia in prossimità di Milano, sia più in generale in Italia. 

Infine, riprendendo alcune delle analisi realizzate su dati Istat, si sono ripetuti i 
calcoli incentrati, sull'evolversi del numero di sinistri al 
cambiare dell'orario, o del mese dell'anno.

\skipline
I dataset riguardanti le informazioni meteo, sono stati utilizzati per un'analisi sulla 
correlazione tra i fattori atmosferici e l'incidentalità per le strade in 
prossimità della città di Milano. 
In particolare, ci si è concentrati sui dati riguardanti la temperatura, 
l'umidità e la velocità del vento. 

Causa la bassa precisione dei dataset utilizzati, tuttavia, questo capitolo è 
da intendersi come una dimostrazione di quanto, un calcolo realizzato tra due campioni 
scollegati, possa fruttare un risultato plausibile ma comunque infondato. 

D'altra parte, non è da escludere che, tramite analisi più approfondite, 
con informazioni più precise, sia possibile ottenere risultati validi.

\section{Conclusioni}

\`E opportuno evidenziare che la maggior parte delle analisi realizzate, 
è stata portata a termine scendendo a compromessi tra la disponibilità del dato e la 
precisione di questo. 

Se per analisi realizzate su dataset Istat, i campi sono già 
parzialmente aggregati, e quindi di facile elaborazione, 
per i dati riguardanti le coordinate degli incidenti, 
o quelli sugli autovelox, uno dei maggiori problemi è stata l'assenza delle 
informazioni necessarie come, ad esempio, le date del 
posizionamento delle telecamere tutor a Milano. 

Un altro problema, questa volta presente prevalentemente nei calcoli su dati già aggregati, 
è stata la difficoltà nel creare un contesto valido attorno alle informazioni ricavate. 
\`E il caso dei sinistri divisi in base alla fascia di età del conducente, dove è possibile 
ottenere esiti marcatamente differenti a seconda del denominatore utilizzato per 
stimare la popolazione per età.

\raggedleft\vfill\scriptsize Ultima revisione: 5 febbraio 2021\par

\end{document}